\chapter{Організаційне управління та стратегія протидії злоякісному коду}
\chaptershort{Управління та стратегія}
\label{chap:ch17}

Технічні засоби (розділи~\ref{chap:ch06}--\ref{chap:ch16}) дають ефект лише в межах зрілої системи \term{управління} (governance): без стратегії, процесів, ролей, бюджету, людей і культури найкращі технології реалізують мізерну частку потенціалу. Цей завершальний розділ піднімається з технічного на організаційно-стратегічний рівень: він розглядає протидію ЗПК як керовану функцію --- із цілями, ризиками, інвестиціями, метриками й безперервним удосконаленням. Він узагальнює наскрізні принципи посібника (ешелонування, каскад вартості й покриття, перенесення на стійкі ознаки, припущення компрометації) у цілісну стратегію й показує, що протидія ЗПК --- не набір засобів, а дисципліна управління ризиком в умовах активного супротивника. Розділ розглядає управління безпекою й рамкові моделі, управління ризиком та економіку безпеки, чинник людини й культуру, стратегію кіберстійкості, метрики й зрілість, а також національний вимір в Україні, і завершується кейсом побудови програми протидії ЗПК.

\section{Управління безпекою й рамкові моделі}\label{sec:ch17-governance}

\term{Управління безпекою} визначає, \emph{хто}, \emph{що} і \emph{навіщо} робить для захисту: цілі, політики, ролі, відповідальність, процеси, контроль. Рамкові моделі структурують цю діяльність. NIST Cybersecurity Framework~\cite{ch17:nist-csf} виділяє функції --- \emph{управляти}, \emph{ідентифікувати}, \emph{захищати}, \emph{виявляти}, \emph{реагувати}, \emph{відновлювати}, --- що охоплюють повний цикл і прямо відповідають структурі цього посібника (моделювання загроз, профілактика, виявлення, реагування). ISO/IEC~27001~\cite{ch17:iso27001} задає вимоги до системи управління інформаційною безпекою (СУІБ) із циклом постійного вдосконалення (плануй-виконуй-перевіряй-дій).

\begin{tikzfig}
\begin{tikzpicture}[font=\scriptsize,
  b/.style={draw,rounded corners=2pt,minimum height=5.5mm,text width=16mm,align=center,font=\tiny},
  >={Stealth[length=1.4mm]}]
\node[b] (g) at (0,0) {управляти};
\node[b] (i) at (1.9,0) {ідентифік.};
\node[b] (p) at (3.8,0) {захищати};
\node[b] (d) at (5.7,0) {виявляти};
\node[b] (r) at (7.6,0) {реагувати};
\node[b] (rec) at (3.8,-1.3) {відновлювати};
\draw[->] (i)--(p); \draw[->] (p)--(d); \draw[->] (d)--(r); \draw[->] (r) to[out=-90,in=0] (rec);
\draw[->] (rec) to[out=180,in=-90] (i);
\draw[-,dashed] (g) -- (i);
\node[font=\tiny,align=center] at (0,-1.0) {наскрізна\\функція};
\end{tikzpicture}
\figcap{Рис.~\thechapter.1 --- Функції управління кібербезпекою (за NIST CSF) як безперервний цикл}
\end{tikzfig}

Значення рамкових моделей --- не в бюрократії, а в \emph{повноті} й \emph{узгодженості}: вони гарантують, що жодну функцію не пропущено (наприклад, організація не зосередилася лише на запобіганні, занедбавши виявлення й реагування), і що діяльність узгоджена з цілями й ризиком. Для аспіранта-дослідника вони дають спільну мову для систематизації технічних методів у цілісну систему. Основні рамкові моделі та їх призначення зіставлено в табл.~\thechapter.1.

\tabhead{Таблиця~\thechapter.1 --- Рамкові моделі управління кібербезпекою}

\begin{center}\footnotesize
\begin{tabularx}{\linewidth}{L{22mm} Y L{18mm}}
\toprule
Модель & Призначення & Форма \\
\midrule
NIST CSF & функції захисту, спільна мова & добровільна рамка \\
ISO/IEC~27001 & система управління (СУІБ) & сертифікована \\
C2M2 & оцінка й розвиток зрілості & самооцінка \\
\bottomrule
\end{tabularx}
\end{center}

\section{Управління ризиком та економіка безпеки}\label{sec:ch17-risk}

Ресурси завжди обмежені, тому протидія ЗПК --- це управління \emph{ризиком} (розділ~\ref{chap:ch04}), а не прагнення недосяжної абсолютної безпеки. Рішення про інвестиції в захист ґрунтуються на оцінці ризику (ймовірність $\times$ наслідки) і співвідношенні вартості заходу з очікуваним зниженням ризику~\cite{ch17:fair}. Економіка інформаційної безпеки~\cite{ch17:anderson} показує, що багато проблем мають не технічну, а економічну природу: хибні стимули, перенесення ризику, зовнішні ефекти.

Класична модель Гордона--Лоба~\cite{ch17:gordon-loeb} формалізує оптимальний рівень інвестицій. Якщо $v$ --- потенційні збитки, а інвестиція $z$ знижує ймовірність реалізації загрози, то чиста вигода

\begin{equation}
\mathrm{ENBIS}(z) = v\bigl[S(0) - S(z)\bigr] - z,
\label{eq:ch17-gl}
\end{equation}

де $S(z)$ --- залишкова ймовірність успішної атаки за інвестиції $z$, максимізується за скінченного рівня інвестицій, що \emph{менший} за очікувані збитки. Формула~\eqref{eq:ch17-gl} формалізує фундаментальну істину: раціональна мета --- не нульовий ризик (невигідно й недосяжно), а оптимальний баланс, де гранична вигода від додаткового захисту дорівнює його граничній вартості. Це пояснює, чому розумна стратегія пріоритезує заходи з найбільшим відношенням зниження ризику до вартості --- зокрема дешеву профілактику й ешелонування (розділ~\ref{chap:ch10}) перед дорогими точковими рішеннями.

\section{Чинник людини й культура безпеки}\label{sec:ch17-people}

Людина --- водночас найслабша ланка й найцінніший актив захисту. Значна частина інцидентів починається з людини: фішинг, слабкі паролі, помилки конфігурації, зневага до оновлень~\cite{ch17:verizon}. Жодна технологія не компенсує повністю необережного користувача чи адміністратора. Тому \term{культура безпеки} --- спільне усвідомлення загроз і відповідальна поведінка --- є невіддільною частиною протидії.

Проте підвищення обізнаності часто неефективне~\cite{ch17:awareness}: разові тренінги не змінюють поведінки. Дієві підходи: вбудовування безпеки в процеси (щоб безпечна дія була й найзручнішою), реалістичні навчання, позитивна культура (без звинувачення, що заохочує повідомляти про помилки й підозри), узгодження стимулів. Ключовий зсув --- від погляду на людину як на проблему до погляду на неї як на \emph{сенсор} і учасника захисту: навчений працівник, що вчасно повідомляє про підозрілий лист, --- цінний елемент виявлення (розділ~\ref{chap:ch11}).

Окрему категорію становлять \emph{внутрішні} загрози --- зловживання з боку осіб із легітимним доступом (незадоволений працівник, завербований інсайдер, скомпрометований обліковий запис адміністратора). Вони особливо небезпечні, бо діють із-за межі довіри, обходячи периметровий захист, і зливаються з легітимною діяльністю --- та сама проблема контексту, що й у LotL (розділ~\ref{chap:ch11}). Протидія поєднує технічне (найменші привілеї, розділення обов'язків, моніторинг аномалій поведінки користувачів) з організаційним (перевірка при прийомі, чіткі політики, здорова культура, що знижує мотивацію до зловживань). Внутрішні загрози показово ілюструють, що управління безпекою не зводиться до технологій: воно невіддільне від управління персоналом, процесами й організаційним кліматом.

Зрештою, роль \emph{керівництва} є вирішальною. Протидія ЗПК потребує бюджету, повноважень і пріоритету, які може надати лише вище керівництво; без його залучення програма безпеки лишається технічною ініціативою без ресурсів і впливу. Тому зрілі рамкові моделі~\cite{ch17:nist-csf} виділяють управління (governance) в окрему наскрізну функцію: безпека починається не з інструментів, а з рішення керівництва ставитися до кіберризику як до одного з ключових бізнес-ризиків, нарівні з фінансовим чи репутаційним.

\section{Стратегія кіберстійкості}\label{sec:ch17-resilience}

Сучасна стратегія виходить із \term{припущення компрометації} (assume breach, розділ~\ref{chap:ch13}): абсолютне запобігання недосяжне, тому мета зміщується від недопущення будь-якої компрометації до \term{кіберстійкості}~\cite{ch17:resilience} --- здатності витримувати атаки, обмежувати їх вплив, швидко відновлюватися й продовжувати виконувати місію навіть під час атаки. Це фундаментальний зсув парадигми: не фортеця, яку або взяли, або ні, а організм, що виживає й адаптується попри ушкодження.

\begin{tikzfig}
\begin{tikzpicture}[font=\scriptsize,
  b/.style={draw,rounded corners=2pt,minimum height=5.5mm,text width=18mm,align=center,font=\tiny},
  >={Stealth[length=1.4mm]}]
\node[b] (prev) at (0,0) {запобігати (знизити ймовірність)};
\node[b] (with) at (2.6,0) {витримувати (обмежити вплив)};
\node[b] (rec) at (5.2,0) {відновлювати (швидко)};
\node[b] (adapt) at (7.8,0) {адаптувати (вчитися)};
\draw[->] (prev)--(with); \draw[->] (with)--(rec); \draw[->] (rec)--(adapt);
\draw[->] (adapt) to[out=-90,in=-90] (prev.south);
\end{tikzpicture}
\figcap{Рис.~\thechapter.2 --- Кіберстійкість: від запобігання до здатності витримувати, відновлюватися й адаптуватися}
\end{tikzfig}

Принципи стійкості: ешелонування (розділ~\ref{chap:ch10}), сегментація (обмеження поширення), надмірність і резервування, швидке виявлення й реагування (розділ~\ref{chap:ch16}), припущення компрометації й нульова довіра (розділ~\ref{chap:ch15}), безперервне навчання на інцидентах. Стійкість поєднує всі технічні розділи посібника в єдину стратегію: кожен метод --- рубіж, що або знижує ймовірність, або обмежує вплив, або пришвидшує відновлення.

Практичний наслідок парадигми стійкості --- перенесення частини інвестицій із запобігання на \emph{обмеження наслідків} і \emph{відновлення}. Організація, що вклала все в периметр, але не здатна швидко виявити, стримати й відновитися, зазнає катастрофічних збитків від єдиного прориву --- саме це сталося з багатьма жертвами NotPetya (розділ~\ref{chap:ch15}), чиї пласкі мережі дозволили коду поширитися без перешкод, а відсутність перевірених офлайн-копій унеможливила швидке відновлення. Натомість стійка організація проєктує системи так, щоб компрометація частини не означала втрати цілого: сегментація локалізує ураження, надмірність зберігає функцію, перевірені резервні копії уможливлюють відновлення. Ця зміна акценту --- від «не впустити» до «вистояти й відновитися» --- є найважливішим стратегічним уроком сучасної протидії ЗПК.

\section{Метрики, зрілість і національний вимір}\label{sec:ch17-maturity}

Управління потребує \term{вимірювання}: не можна керувати тим, чого не міряють. Метрики протидії ЗПК охоплюють технічні (повнота й точність виявлення, час присутності~\eqref{eq:ch10-dwell}, покриття технік) та управлінські (зрілість процесів, частка закритих вразливостей, час реагування, результати навчань). \term{Моделі зрілості}~\cite{ch17:c2m2} описують еволюцію від реактивного й хаотичного стану до керованого, вимірюваного й оптимізованого, даючи дорожню карту вдосконалення. Приклади метрик обох груп наведено в табл.~\thechapter.2.

\tabhead{Таблиця~\thechapter.2 --- Метрики протидії ЗПК}

\begin{center}\footnotesize
\begin{tabularx}{\linewidth}{L{22mm} Y}
\toprule
Група & Приклади \\
\midrule
Технічні & повнота й точність виявлення, час присутності, покриття технік ATT\&CK \\
Операційні & час реагування, потік тривог на аналітика, PPV \\
Управлінські & зрілість процесів, частка закритих вразливостей, результати навчань \\
\bottomrule
\end{tabularx}
\end{center}

Важливо добирати метрики обачно: погано дібрана метрика спотворює поведінку (що міряють, те й оптимізують, інколи на шкоду суті). Наприклад, гонитва за кількістю закритих тривог заохочує поверхневе їх закриття, а не якісне розслідування. Тому зрілі програми вимірюють \emph{результати} (скорочення часу присутності, зниження реалізованого ризику), а не лише \emph{активність} (кількість дій), і регулярно переглядають самі метрики.

Зрілість зручно характеризувати як ступінь просування рівнями. Якщо виділити $L$ рівнів зрілості, а організація досягла рівня $\ell$ з якістю реалізації $\rho\in[0;1]$, то інтегральний показник зрілості

\begin{equation}
M = \frac{\ell - 1 + \rho}{L - 1} \in [0;1],
\label{eq:ch17-maturity}
\end{equation}

дозволяє відстежувати прогрес і порівнювати стан у часі. Формула~\eqref{eq:ch17-maturity} підкреслює, що зрілість --- не разове досягнення рівня, а безперервний рух, бо загрози еволюціонують, і стан, достатній сьогодні, застаріває завтра.

Національний вимір для України має стратегічне значення: Стратегія кібербезпеки України~\cite{ch17:strategy-ua} визначає державну рамку протидії, а найінтенсивніший у світі досвід відбиття атак на критичну інфраструктуру робить українські організаційно-управлінські практики~\cite{ch17:governance-ua} цінними й для світу. Протидія ЗПК на національному рівні поєднує державне регулювання (розділ~\ref{chap:ch02}), розвідку загроз (CERT-UA, розділ~\ref{chap:ch12}), захист критичної інфраструктури (розділ~\ref{chap:ch14}) й міжнародну співпрацю в єдину систему кіберстійкості держави.

\section{Кейс: побудова програми протидії ЗПК}\label{sec:ch17-case}

Розглянемо побудову організаційної програми протидії ЗПК (управлінська задача, що синтезує весь посібник).

\emph{Крок 1. Контекст і ризик.} Визначаємо місію, активи, модель загроз (розділ~\ref{chap:ch04}), регуляторні вимоги (розділ~\ref{chap:ch02}) й апетит до ризику. Пріоритезуємо за ризиком, а не за технічною модою.

\emph{Крок 2. Рамка й цілі.} Обираємо рамкову модель (NIST CSF, ISO~27001), ставимо вимірювані цілі за всіма функціями: ідентифікація, захист, виявлення, реагування, відновлення.

\emph{Крок 3. Ешелонований захист.} Будуємо рубежі за каскадом вартості й покриття: профілактика й гартування (розділ~\ref{chap:ch10}), виявлення (сигнатури, поведінка, МН --- розділи~\ref{chap:ch06}--\ref{chap:ch11}), розвідка й полювання (розділи~\ref{chap:ch12}--\ref{chap:ch13}), реагування й форензика (розділ~\ref{chap:ch16}), пріоритезуючи заходи з найкращим відношенням зниження ризику до вартості~\eqref{eq:ch17-gl}.

\emph{Крок 4. Люди й процеси.} Визначаємо ролі (SOC/CSIRT), процеси реагування, навчання, культуру безпеки; перетворюємо людей на сенсори захисту.

\emph{Крок 5. Вимірювання й удосконалення.} Запроваджуємо метрики й модель зрілості~\eqref{eq:ch17-maturity}, безперервно вимірюємо, вчимося на інцидентах і вдосконалюємося, адаптуючись до еволюції загроз.

Кейс синтезує весь посібник: ефективна протидія ЗПК --- це не набір засобів, а керована, вимірювана, безперервно вдосконалювана програма, що балансує ризик і вартість, поєднує технології, процеси й людей і прагне не абсолютної безпеки, а кіберстійкості в умовах активного супротивника.

\section*{Висновки до розділу}
\addcontentsline{toc}{section}{Висновки до розділу}

Технічні засоби реалізують свій потенціал лише в межах зрілого управління: протидія ЗПК --- це керована функція з цілями, ризиками, інвестиціями, метриками й безперервним удосконаленням, а не набір технологій. Рамкові моделі (NIST CSF, ISO/IEC~27001) забезпечують повноту й узгодженість функцій --- від ідентифікації до відновлення. Протидія є управлінням ризиком, а не прагненням абсолютної безпеки: економіка безпеки~\eqref{eq:ch17-gl} показує, що раціональна мета --- оптимальний баланс вартості й зниження ризику, а не нульовий ризик. Чинник людини --- найслабша ланка й цінний сенсор водночас --- вимагає культури безпеки, вбудованої в процеси, а не разових тренінгів. Сучасна стратегія виходить із припущення компрометації й зміщується до кіберстійкості --- здатності витримувати, відновлюватися й адаптуватися, --- що поєднує всі технічні розділи в єдине ціле. Управління потребує вимірювання й моделей зрілості~\eqref{eq:ch17-maturity} для безперервного вдосконалення. Для України протидія ЗПК має стратегічний національний вимір. Головний і наскрізний висновок посібника: ефективна протидія злоякісному коду --- це дисципліна управління ризиком в умовах активного, мислячого супротивника, що поєднує технології, процеси й людей у кероване, вимірюване, кіберстійке ціле, яке не прагне недосяжної абсолютної безпеки, а невпинно підвищує вартість атаки й здатність організації вистояти та відновитися.

\selfcheck
\begin{enumerate}
\item Чому технічні засоби реалізують потенціал лише в межах зрілого управління?
\item Що визначає управління безпекою? Яку роль відіграють рамкові моделі (NIST CSF, ISO~27001)?
\item Як функції NIST CSF (рис.~\thechapter.1) відповідають структурі протидії ЗПК?
\item Чому протидія ЗПК є управлінням ризиком, а не прагненням абсолютної безпеки?
\item Поясніть модель Гордона--Лоба~\eqref{eq:ch17-gl}. Чому оптимальні інвестиції менші за очікувані збитки?
\item Чому багато проблем безпеки мають економічну, а не технічну природу?
\item Чому людина є водночас найслабшою ланкою й цінним активом? Чому разові тренінги неефективні?
\item Які підходи дієво формують культуру безпеки? Що означає погляд на людину як на сенсор?
\item Що таке припущення компрометації? Чим кіберстійкість відрізняється від парадигми фортеці?
\item Назвіть принципи кіберстійкості (рис.~\thechapter.2). Як вони поєднують технічні розділи посібника?
\item Навіщо потрібні метрики й моделі зрілості? Поясніть показник зрілості~\eqref{eq:ch17-maturity}.
\item У чому національний вимір протидії ЗПК для України? Як поєднуються його складові?
\item Опишіть побудову програми протидії ЗПК (кейс підрозділу~\ref{sec:ch17-case}). Як він синтезує весь посібник?
\item Дослідницьке: розробіть програму протидії ЗПК для модельної організації за обраною рамковою моделлю й оцініть її зрілість~\eqref{eq:ch17-maturity} та пріоритети інвестицій~\eqref{eq:ch17-gl}.
\item Дослідницьке: проаналізуйте співвідношення вартості та зниження ризику для набору заходів протидії ЗПК і побудуйте оптимальний за Гордоном--Лобом портфель інвестицій.
\end{enumerate}
